\documentclass[11pt]{article}
\usepackage[a4paper,margin=1in]{geometry}
\usepackage{amsmath,amssymb,mathtools,bm}
\usepackage{enumitem,booktabs,array}
\usepackage{tcolorbox}
\usepackage{hyperref}
\usepackage[protrusion=true,expansion=false]{microtype}
\hypersetup{colorlinks=true,linkcolor=blue,urlcolor=blue,citecolor=blue}
\setlist[itemize]{topsep=2pt,itemsep=2pt}
\setlist[enumerate]{topsep=2pt,itemsep=2pt}
\newtcolorbox{keybox}{colback=blue!3!white,colframe=blue!40!black,boxrule=0.4pt,arc=1mm}
\newtcolorbox{alertbox}{colback=red!3!white,colframe=red!40!black,boxrule=0.4pt,arc=1mm}
\newtcolorbox{answerbox}{colback=green!3!white,colframe=green!40!black,boxrule=0.4pt,arc=1mm}
\newtcolorbox{drillbox}{colback=yellow!5!white,colframe=orange!60!black,boxrule=0.4pt,arc=1mm}
\newcommand{\EV}{\mathbb{E}}
\newcommand{\Var}{\mathrm{Var}}
\newcommand{\Cov}{\mathrm{Cov}}
\newcommand{\blank}[1]{\underline{\hspace{#1}}}

\title{E06-01: Hoberg \& Phillips (2010, RFS) \\
\large ``Product Market Synergies and Competition in Mergers and Acquisitions: A Text-Based Analysis'' --- Active Recall Workbook}
\author{Empirical Corporate Finance II --- Exam Prep}
\date{\today}
\begin{document}\maketitle

\begin{keybox}
\textbf{Paper in one sentence.} HP build text-based product-market similarity scores from 10-K business descriptions and show that M\&A targets are systematically closer in product-space to acquirers than randomly-matched control firms, and that closer product-market similarity predicts larger post-merger cash-flow gains --- providing text-based evidence that product-market synergies are a real driver of M\&A value creation.

\textbf{Placement.} Foundational paper in text-based industry classification; precursor to Hoberg-Phillips Text-based Network Industry Classification (TNIC). Key reference for M\&A-synergy measurement and product-market research.
\end{keybox}

\section*{Round 1: Complete Walk-Through}

\subsection*{1.1 Research question}
Standard M\&A studies rely on SIC/NAICS codes, which are coarse and time-invariant. Can we measure product-market closeness directly from firm disclosures, and does product-market similarity predict M\&A selection and gains?

\subsection*{1.2 Text-based similarity}
\begin{itemize}
\item Parse Item 1 (business description) of every 10-K.
\item Extract nouns/noun phrases; convert to binary/tf-idf vectors.
\item Cosine similarity $s_{ij}=\frac{V_i\cdot V_j}{\|V_i\|\|V_j\|}$ across firm pairs.
\item Values $\in[0,1]$; TNIC industry = set of $j$ with $s_{ij}$ above threshold.
\end{itemize}

\subsection*{1.3 Data and design}
\begin{itemize}
\item All 10-K filings 1996--2005.
\item SDC M\&A sample in the same window.
\item Logit: probability target was acquired as function of similarity with bidder.
\item Cross-sectional regressions of announcement returns and long-run operating performance on similarity.
\end{itemize}

\subsection*{1.4 Headline results}
\begin{itemize}
\item Targets are much more similar to acquirers than random control firms (conditional on same SIC industry even, within-industry similarity varies a lot).
\item Higher product-market similarity predicts larger combined-firm CARs, higher post-merger cash flows, and greater synergies.
\item Similarity captures synergy variation missed by SIC-based measures.
\item Effects robust to fixed effects and measurement choices.
\end{itemize}

\subsection*{1.5 Mechanism}
\begin{itemize}
\item Similar product lines imply economies of scale, scope, complementarity.
\item Text-based measure captures evolving product markets that SIC codes cannot.
\item Consistent with a synergy-based explanation of M\&A rather than pure empire-building.
\end{itemize}

\subsection*{1.6 Implications}
\begin{itemize}
\item Text-based industry classification is a richer alternative to SIC/NAICS.
\item Allows research on product differentiation, competition, and synergy at a fine level.
\item TNIC data (Hoberg-Phillips website) now standard in applied empirical finance.
\end{itemize}

\subsection*{1.7 Caveats}
\begin{alertbox}
\begin{itemize}
\item 10-K business descriptions are strategic disclosures, potentially selectively written.
\item Similarity measured only among public filers.
\item Text-based does not capture production-cost or geographic dimensions of markets.
\end{itemize}
\end{alertbox}

\section*{Round 2: Level 1 Fill-in}
\begin{enumerate}
\item Text source: Item 1 of firm \blank{1cm}-K.
\item Similarity measure: \blank{2cm} of word vectors.
\item Period: \blank{1cm}--\blank{1cm}.
\item Finding: targets more similar to acquirers than \blank{2cm} controls.
\item Outcome: similarity predicts higher combined \blank{2cm} and cash flows.
\end{enumerate}

\section*{Round 3: Level 2 Prompts}
\begin{enumerate}
\item Describe the text-based similarity construction.
\item Report the main findings on selection and post-merger performance.
\item Compare text-based to SIC/NAICS industry classifications.
\item Discuss the mechanism and implications for M\&A research.
\item What are the caveats, and how does HP extend to TNIC?
\end{enumerate}

\section*{Round 4: Minimal Scaffolding}
From a blank page: (i) motivation, (ii) text method, (iii) findings, (iv) mechanism, (v) extensions.

\section*{Round 5: Blank Page}
Essay: text-based product-market similarity and M\&A synergy.

\vspace{6pt}
\begin{drillbox}
\textbf{Speed Drill.} (1) Text source. (2) Similarity metric. (3) Sample years. (4) Selection finding. (5) Synergy finding.
\end{drillbox}

\section*{Quick Reference \& Answer Key}
\begin{answerbox}
\begin{itemize}
\item \textbf{Text.} 10-K Item 1 business descriptions.
\item \textbf{Measure.} Cosine similarity of word vectors.
\item \textbf{Selection.} Targets closer to acquirers.
\item \textbf{Synergy.} Similarity $\to$ higher CAR \& CF.
\end{itemize}
\end{answerbox}

\begin{keybox}
\textbf{30-second exam pitch.} Hoberg \& Phillips (2010) build text-based product-market similarity scores from Item 1 of 10-Ks, 1996--2005. Cosine similarity of noun-vector representations identifies firm pairs that are close in product space, even within SIC industries. Applied to SDC mergers, targets are dramatically more similar to acquirers than randomly-matched control firms, and similarity strongly predicts higher combined-firm CARs and post-merger cash flows, indicating real product-market synergies. This challenges the sufficiency of SIC/NAICS for synergy research and establishes text-based industry classification as a richer alternative. The paper seeds the Hoberg-Phillips TNIC dataset now standard in applied empirical finance, and provides foundational evidence that M\&A creates product-market value consistent with synergy theories rather than pure empire-building.
\end{keybox}

\end{document}
